\section{Scope and Limitations}
\label{sec:limitations}

These limitations define the scope of the public-record claims rather than weakening the main result. The repository collection is a May 2026 \hf snapshot scoped to compact and derived open \llm artifacts, not a census of every \hf repository. The paper corpus is a constructed scholarly record set; main percentages use 2,214 analysis-set paper records unless otherwise stated. Artifact role and package form are audited inference fields derived from public metadata, with stronger support for fine-tune, adapter/LoRA, and quantized/deployment records than for base and merge/mixed records. Provenance tiers measure public upstream evidence; they do not recover private training logs, hidden checkpoints, or developer knowledge. Release-package evidence records visible assets, but execution, environment reconstruction, model quality, adoption, and preservation success remain downstream tasks.

\section{Conclusion}
\label{sec:conclusion}

This paper contributes a record-design measurement framework for compact and derived open \llm artifacts. These artifacts are becoming a distributed digital collection represented across model hubs, model cards, scholarly papers, code repositories, family mentions, release statements, and paper--repository links. Linking a fixed \hf repository snapshot with scholarly paper records shows that broad visibility is common, but coordinated evidence is much rarer: 90.7\% of paper records are discovery-ready, 18.1\% meet the operational P+R threshold, 6.1\% satisfy the high-curatability threshold, and 2.3\% qualify as preservation-review candidates.

The implication for digital libraries is concrete. Open-model infrastructure should expose typed artifact roles, package-form fields, upstream relation types, provenance-evidence tiers, bidirectional paper--repository links, and release-package fields. These interventions shift the question from whether open \llm artifacts are online to whether public records can justify curation decisions about attribution, traceability, reuse assessment, and preservation triage. The next generation of open-model infrastructure should not only keep artifacts online; it should keep the evidence around those artifacts connected.
